\documentclass[../../main.tex]{subfiles}
\graphicspath{{../../image/}{../../figure/}}
\begin{document}

\section{自相似解与 Riemann 问题}\label{section:ch09-自相似解与riemann问题}

双曲守恒律方程组
\begin{align}
\partial_{t}U+\partial_{x}F(U)=0
\label{align:ch09-9-1-1}
\end{align}
在坐标均匀拉伸 $(x,t)\mapsto(\alpha x,\alpha t)$ 下不变;因而它容许\emph{自相似解},其定义在时空平面上,并沿从原点发出的直线射线保持常值。又因\eqref{align:ch09-9-1-1}在坐标平移 $(x,t)\mapsto(x+\bar{x},t+\bar{t})$ 下也不变,自相似解的焦点可从原点平移到时空中任一固定点 $(\bar{x},\bar{t})$。

\eqref{align:ch09-9-1-1} 的一个(一般为弱的)自相似解 $U$,定义在上半平面或下半平面上、聚焦于原点时,有形如
\begin{align}
U(x,t)=V\!\left(\frac{x}{t}\right),\qquad-\infty<x<\infty,\quad 0<t<\infty\ \text{或}\ -\infty<t<0,
\label{align:ch09-9-1-2}
\end{align}
的表示,其中 $V$ 是 $(-\infty,\infty)$ 上的有界可测函数,在分布意义下满足常微分方程
\begin{align}
[F(V(\xi))-\xi V(\xi)]^{\cdot}+V(\xi)=0.
\label{align:ch09-9-1-3}
\end{align}
为验证对上半平面上形如\eqref{align:ch09-9-1-2}的 $U$ 此结论成立,任取支集在 $(-\infty,\infty)\times(0,\infty)$ 内的检验函数 $\phi$,注意到
\begin{align}
\int_{0}^{\infty}\int_{-\infty}^{\infty}\bigl[\partial_{t}\phi(x,t)U(x,t)+\partial_{x}\phi(x,t)F(U(x,t))\bigr]\mathrm{d}x\mathrm{d}t
=\int_{-\infty}^{\infty}\bigl\{\dot{\psi}(\xi)[F(V(\xi))-\xi V(\xi)]-\psi(\xi)V(\xi)\bigr\}\mathrm{d}\xi,
\label{align:ch09-9-1-4}
\end{align}
其中
\begin{align}
\psi(\xi)=\int_{0}^{\infty}\phi(\xi t,t)\,\mathrm{d}t,\qquad-\infty<\xi<\infty.
\label{align:ch09-9-1-5}
\end{align}
定义在下半平面上形如\eqref{align:ch09-9-1-2}的 $U$ 情形类似处理,同样得到\eqref{align:ch09-9-1-3}。

由\eqref{align:ch09-9-1-3}推知 $F(V)-\xi V$ 是 Lipschitz 的,且
\begin{align}
F(V(\xi))-\xi V(\xi)-F(V(\zeta))+\zeta V(\zeta)=-\int_{\zeta}^{\xi}V(\theta)\,\mathrm{d}\theta
\label{align:ch09-9-1-6}
\end{align}
对一切 $\zeta,\xi\in(-\infty,\infty)$ 成立。

\eqref{align:ch09-9-1-3} 的任一有界可测解 $V$ 的定义域 $(-\infty,\infty)$ 被划分为: $V$ 的间断点集 $\mathscr{S}$; $V$ 的连续点集 $\mathscr{W}$,即满足
\begin{align}
\lambda_{i}(V(\xi))=\xi
\label{align:ch09-9-1-7}
\end{align}
(对某个 $i\in\{1,\dots,n\}$) 的连续点 $\xi$;以及 $V$ 的连续点集 $\mathscr{C}$,即对 $i=1,\dots,n$ 均有 $\lambda_{i}(V(\xi))\neq\xi$ 的连续点 $\xi$。

任一收敛到 $\mathscr{S}$ 中一点 $\xi$ 的递增序列 $\{\zeta_{k}\}$ 与递减序列 $\{\xi_{k}\}$ 都含有子列(仍记作 $\{\zeta_{k}\}$ 与 $\{\xi_{k}\}$),使得 $\{V(\zeta_{k})\}$ 与 $\{V(\xi_{k})\}$ 分别收敛到极限 $V_{-}$ 与 $V_{+}$。由\eqref{align:ch09-9-1-6},
\begin{align}
F(V_{+})-F(V_{-})=\xi[V_{+}-V_{-}].
\label{align:ch09-9-1-8}
\end{align}
特别地,当单侧极限 $V(\xi\pm)$ 存在时, $\xi$ 标示一个激波,而\eqref{align:ch09-9-1-8}正是 Rankine-Hugoniot 跳跃条件。然而在目前这种一般性程度下,我们不能排除收敛到 $\xi$ 的不同序列可能产生不同的极限值 $V_{-}$ 和/或 $V_{+}$。无论如何,由\eqref{align:ch09-9-1-8}可知,只要 $|V_{+}-V_{-}|$ 充分小, $\xi$ 必接近于某个 $i\in\{1,\dots,n\}$ 的 $\lambda_{i}(V_{\pm})$。这又推出: $\mathscr{S}$ 闭包中 $V$ 的任一连续点必属于 $\mathscr{W}$。同时, $\mathscr{W}$ 闭包中 $V$ 的任一连续点也属于 $\mathscr{W}$。因此 $\mathscr{S}\cup\mathscr{W}$ 是闭集,而 $\mathscr{C}$ 是开集。

作为开集, $\mathscr{C}$ 是至多可数个互不相交开区间的并。任取其中一个开区间,记为 $\mathscr{I}$,并固定 $\xi,\zeta\in\mathscr{I}$。我们把\eqref{align:ch09-9-1-6}改写为
\begin{align}
(A-\xi I)[V(\xi)-V(\zeta)]=-\int_{\zeta}^{\xi}[V(\theta)-V(\zeta)]\,\mathrm{d}\theta,
\label{align:ch09-9-1-9}
\end{align}
其中 $A$ 是 $n\times n$ 矩阵
\begin{align}
A=\int_{0}^{1}\bigl[s\,\mathrm{D}F(V(\zeta))+(1-s)\mathrm{D}F(V(\xi))\bigr]\mathrm{d}s.
\label{align:ch09-9-1-10}
\end{align}
当 $\zeta\to\xi$ 时, $\xi$ 与 $A$ 的特征值保持有正距离,而\eqref{align:ch09-9-1-9}右端为 $o(1)|\xi-\zeta|$。故\eqref{align:ch09-9-1-9}推出 $V$ 在 $\mathscr{I}$ 上 Lipschitz,且其导数几乎处处为零,即在 $\mathscr{I}$ 上 $V$ 为常值。

上述约束仍为 $V$ 的结构留出了充分的复杂余地,这取决于方程组\eqref{align:ch09-9-1-1}的性质。然而,当\eqref{align:ch09-9-1-1}严格双曲且 $V$ 的值域被限制在以某状态 $\bar{U}$ 为中心、半径 $\delta$ 很小的球 $\mathscr{B}_{\delta}(\bar{U})$ 内时,就会出现简单而规整的构型。$\delta$ 的大小受与\cref{section:ch08-强激波弱激波中等强度激波}中刻画中等强度激波同类约束的限制。在这些假设下, $\mathscr{S}$ 与 $\mathscr{W}$ 被限制在以 $\lambda_{1}(\bar{U}),\dots,\lambda_{n}(\bar{U})$ 为中心、长度为 $O(\delta)$ 的 $n$ 个区间之并内。我们以 $\mathscr{S}_{i}$、$\mathscr{W}_{i}$ 记 $\mathscr{S}$、$\mathscr{W}$ 中位于 $\lambda_{i}(\bar{U})$ 邻域内的部分,并取闭集 $\mathscr{S}_{i}\cup\mathscr{W}_{i}$ 的凸包 $[\zeta_{i},\xi_{i}]$。于是 $V$ 在区间 $(-\infty,\zeta_{1}),(\xi_{1},\zeta_{2}),\dots,(\xi_{n},\infty)$ 上取常值 $U_{0},U_{1},\dots,U_{n}$。区间 $[\zeta_{i},\xi_{i}]$ 内还可能含有 $\mathscr{C}$ 的一个开子集,记作 $\mathscr{C}_{i}$。我们称 $[\zeta_{i},\xi_{i}]$ 为一个\emph{$i$-波}。结论是:在上述假设下, $V$ 由 $n$ 个波组成一个波扇,每个特征族各一个,以近似特征速度传播,并由常值状态分隔。当自相似解定义在上半平面时,该波扇自原点向外发出;当自相似解定义在下半平面时,该波扇向内汇入原点。

当函数 $V$ 有界变差时, $i$-波的结构尤为清晰。此时 $\mathscr{S}_{i}$ 可数,且在每个 $\xi\in\mathscr{S}_{i}$ 处单侧极限 $V_{\pm}=V(\xi\pm)$ 存在,满足跳跃条件\eqref{align:ch09-9-1-8}。故 $\mathscr{S}_{i}$ 由 $i$-激波组成。此外, $V$ 几乎处处可微。特别地,若 $V$ 在点 $\xi\in\mathscr{W}_{i}$ 处可微,则\eqref{align:ch09-9-1-3}与\eqref{align:ch09-9-1-7}给出
\begin{align}
[\mathrm{D}F(V(\xi))-\xi I]\dot{V}(\xi)=0,
\label{align:ch09-9-1-11}
\end{align}
\begin{align}
\mathrm{D}\lambda_{i}(V(\xi))\dot{V}(\xi)=1.
\label{align:ch09-9-1-12}
\end{align}
因而 $\dot{V}(\xi)$ 是非零向量,与 $\mathrm{D}F(V(\xi))$ 的第 $i$ 个右特征向量 $R_{i}(V(\xi))$ 共线。此外, $V(\xi)$ 必为第 $i$ 特征族真正非线性的状态。据此,我们把 $\mathscr{W}_{i}$ 实现为 $i$-波的中心 $i$-简单波分量:对于向外波扇它是稀疏波,对于向内波扇它是压缩波。

现在我们证明 $V$ 在 $[\zeta_{i},\xi_{i}]$ 上有界变差,至少在该 $i$-波\emph{单向}时如此,即 $DF(\bar{U})$ 的左 $i$-特征向量 $L_{i}(\bar{U})$ 可定向得使
\begin{align}
L_{i}(\bar{U})[V(\xi)-V(\zeta)]\geqslant 0
\label{align:ch09-9-1-13}
\end{align}
对 $V$ 的任一对连续点 $(\zeta,\xi)$\footnote{当 $i=1$ 时 $\xi_{i-1}=-\infty$;当 $i=n$ 时 $\xi_{i+1}=\infty$。}成立,其中 $\xi_{i-1}<\zeta<\xi<\xi_{i+1}$。

固定如上的连续点 $\zeta,\xi$,满足 $\xi_{i-1}<\zeta<\xi<\xi_{i+1}$,由\eqref{align:ch09-9-1-6}经\eqref{align:ch09-9-1-9}得
\begin{align}
[\mathrm{D}F(\bar{U})-\lambda_{i}(\bar{U})I][V(\xi)-V(\zeta)]
=O(\delta)[V(\xi)-V(\zeta)]-\int_{\zeta}^{\xi}[V(\theta)-V(\zeta)]\,\mathrm{d}\theta.
\label{align:ch09-9-1-14}
\end{align}
我们写
\begin{align}
V(\xi)-V(\zeta)=\sum_{k=1}^{n}a_{k}R_{k}(\bar{U}).
\label{align:ch09-9-1-15}
\end{align}
将\eqref{align:ch09-9-1-15}左端乘以 $L_{i}(\bar{U})$,并回顾\eqref{align:ch07-左右向量规范化}与\eqref{align:ch09-9-1-13},推出 $a_{i}\geqslant0$。另一方面,把\eqref{align:ch09-9-1-14}与\eqref{align:ch09-9-1-15}结合得
\begin{align}
\sum_{k\neq i}|a_{k}|\leqslant c\delta a_{i}+c\delta(\xi-\zeta).
\label{align:ch09-9-1-16}
\end{align}
从而
\begin{align}
TV_{[\zeta_{i},\xi_{i}]}V(\cdot)\leqslant L_{i}(\bar{U})[U_{i}-U_{i-1}]+c\delta^{2}.
\label{align:ch09-9-1-17}
\end{align}

以下各节将对波扇施加容许性条件,例如要求任一跳跃间断必须满足\cref{chapter:ch08-容许激波}中讨论的部分或全部激波容许准则,而最低限度必须满足 Lax $E$ 条件。注意 Lax $E$ 条件诱导函数 $\lambda_{i}(V(\xi))$ 在任一向内或向外 $i$-波 $[\zeta_{i},\xi_{i}]$ 上单调。正如我们后文将看到的,在任一严格双曲方程组中,中等强度的容许向外 $i$-波最终是单向的,但向内 $i$-波则未必如此。然而,当第 $i$ 特征族真正非线性时,在 Lax $E$ 条件下即使向内 $i$-波也必单向。理由是:正如\cref{chapter:ch08-容许激波}中所见,当第 $i$ 特征族真正非线性时, $\lambda_{i}$ 沿 $i$ 波曲线单调变化,至少在原点附近如此。

从今以后,我们把讨论限制在有界变差的波扇上。作为有界变差函数, $V$ 是一个 saltus 函数(代表激波)与一个连续函数之和。而连续部分又是一个绝对连续函数与一个纯奇异函数(即所谓 Cantor 部分)之和。已有结论(参考文献见\cref{section:ch09-注记})表明:当方程组真正非线性时,连续部分 Lipschitz,因而特别地 $V$ 是 $SBV$ 函数。这表明真正非线性所起的作用是促成间断或 Lipschitz 连续,而阻碍单纯连续这一"中间地带"。真正非线性的这一特征的其他表现将在后文遇到,最显著的是在\cref{chapter:ch11-一维标量守恒律}中。至关重要的是要理解:由初值在原点引入的跳跃间断是如何被分解为向外波扇的。这正是

\begin{definition}[Riemann 问题]\label{definition:ch09-riemann问题}
确定\eqref{align:ch09-9-1-1}在 $(-\infty,\infty)\times(0,\infty)$ 上的一个自相似(一般为弱)解 $U$,满足初值条件
\begin{align}
U(x,0)=
\begin{cases}
U_{L},&x<0,\\
U_{R},&x>0,
\end{cases}
\label{align:ch09-9-1-18}
\end{align}
其中 $U_{L},U_{R}$ 是 $\mathscr{O}$ 中给定的状态。
\end{definition}

按上文讨论,我们将求形如\eqref{align:ch09-9-1-2}的 Riemann 问题解,其中 $V$ 在 $(-\infty,\infty)$ 上满足常微分方程\eqref{align:ch09-9-1-3},并附边值条件
\begin{align}
V(-\infty)=U_{L},\qquad V(\infty)=U_{R}.
\label{align:ch09-9-1-19}
\end{align}

非唯一性的幽灵再次引出容许性问题,这将是以下各节讨论的主题。

\end{document}
